\begin{frame}[noframenumbering,plain]
    \setcounter{framenumber}{1}
    \maketitle
\end{frame}

\begin{frame}
    \frametitle{Актуальность}

    \begin{block}{}
        \footnotesize
        Компьютерное моделирование --- ключевой инструмент для решения сложных задач гидродинамики, 
        где аналитические решения часто недоступны. 
        Традиционные сеточные методы (FDM, FVM, FEM) имеют ограничения при моделировании:
        \begin{itemize}
            \item свободных поверхностей
            \item сложной геометрии
            \item движущихся границ
            \item многофазных течений
        \end{itemize}
    \end{block}

    \begin{block}{}
        \footnotesize
        Метод сглаженных частиц (SPH) как бессеточный подход эффективен для:
        \begin{itemize}
            \item моделирования волн и их взаимодействия с конструкциями
            \item анализа многофазных течений и задач со сложной геометрией
            \item решения задач в аэродинамике, механике и биологии
        \end{itemize}
    \end{block}
\end{frame}
\note{
    \footnotesize
    Несмотря на преимущества, SPH имеет недостатки:
    \begin{itemize}
        \item Трудности с граничными условиями
        \item Низкая точность при нерегулярном распределении частиц
        \item Неустойчивость при растяжении
    \end{itemize}
}

\begin{frame}
    \frametitle{Цель и задачи работы}
    \begin{itemize}
        \item Цель: разработка комплекса ПО для моделирования и 
        анализа сложных гидродинамических процессов с применением суперкомпьютерных вычислительных технологий
        \item Основные задачи:
        \begin{enumerate}
            \item Изучить научную литературу и существующие реализации SPH
            \item Разработать программный комплекс для численного моделирования методом SPH:
            \begin{enumerate}
                \item Средства подготовки входных данных
                \item Вычислительные модули
                \item Средства визуализации
                \item Средства пост-процессинга
            \end{enumerate}
            \item Провести численные эксперименты и их анализ:
            \begin{enumerate}
                \item Анализ эффективности распараллеливания
                \item Анализ дисперсионного соотношения
                \item Анализ затухания волн на поверхности воды
                \item Трёхмерная задача о разрушении дамбы
            \end{enumerate}
        \end{enumerate}
    \end{itemize}
\end{frame}
\note{
    Проговариваются вслух положения, выносимые на защиту
}
